\makeabstract
    {
    Supersymmetric Quantum Mechanics on a Simplicial Complex
    }
    {
    Joseph Leckie\textsuperscript{1}, Ivan Contreras\textsuperscript{1}, Myeongmin Choi\textsuperscript{1}
    }
    {
    \textsuperscript{1} Amherst College, Amherst, MA
    }
    {
    The path integral formulation of quantum field theory, while successful and insightful, is computationally  difficult to apply even in very simple cases because of its requirement that one takes integrals of infinite measure. In this paper, we extend a ''toy model'' of this formulation which instead uses a simplicial complex as a discrete configuration space. In past works, graph quantum field theory in the free case found supersymmetry to be conserved for all non-trivial graphs, which is a barrier to a realistic model, since we do not see a perfect mass symmetry between bosons and fermions. We expand the graph formulation to an n-dimensional simplicial complex, and then apply potential fields in order to break supersymmetry.
    }
    {
     We define a de Rham Cohomology describing the energy of n-dimensional states. Then we prove the Morse inequalities, which state that the number of n-dimensional critical points of a Morse function is bounded from below by the number of n-dimensional holes in the simplicial complex for any Morse function. A Morse function is a function from the set of simplices to the real numbers that has a clear flow from high to low dimension, where higher dimension simplices tend to have a higher assigned values than lower dimensional ones. Then, we use the de Rham Cohomology to create Hamiltonians for each dimension, which are square positive definite matrices called the n-dimensional Laplacian. Then, we add potential fields to this Hamiltonian and search for 0-energy states. Simplicial complexes for which there are no 0-energy states require supersymmetry breaking.
    }
    {
    We find that supersymmetry is broken for almost all cases of potential fields added to graphs, excluding very specific cases. This makes graph quantum mechanics viable as a descriptive model of the real world, which exhibits supersymmetry breaking. We additionally prove the Morse inequalities on a simplicial complex, validating the supersymmetric approach to a simplicial complex.
    }
    {
    Quantum field theory on a simplicial complex is a valid representation of 3D space mathematically, since supersymmetry can be broken and the Morse inequalities hold. However, more research is required to discern whether it accurately approximates continuum mechanics. 
    }

    \newpage
    